% !PRINTED 650
\weberpage{650}
%<<CONT
\weq{7}{\begin{aligned}k_0+k_1+k_2+k_3 &= 4\,C,\\ k_0-k_1+k_2-k_3 &= 4\,B\sqrt{c},\end{aligned}}
hence
\weq{8}{\begin{aligned}k_0+k_2 &= 2\,(C+B\sqrt{c}),\\ k_1+k_3 &= 2\,(C-B\sqrt{c}),\end{aligned}}
from which by (5) it follows:
\weq{9}{\begin{aligned}k_0 &= C+B\sqrt{c}+\sqrt{b+a\sqrt{c}}\\ k_1 &= C-B\sqrt{c}+\sqrt{b-a\sqrt{c}}\\ k_2 &= C+B\sqrt{c}-\sqrt{b+a\sqrt{c}}\\ k_3 &= C-B\sqrt{c}-\sqrt{b-a\sqrt{c}}.\end{aligned}}
Conversely, it is also easy to show that these quantities
are the roots of a biquadratic cyclic equation.
For if we set, as an abbreviation,
\weq{10}{r = \sqrt{c}, \quad \varrho = \sqrt{b+a\sqrt{c}}, \quad \varrho' = \sqrt{b-a\sqrt{c}},}
and hence by (6)
\[ r = \varrho\varrho', \quad \varrho^2 = b+ar, \quad \varrho'^2 = b-ar, \]
then we can represent every rational function of $k_0, k_1, k_2, k_3$ as a linear
function with rational coefficients of the six radicals
\[ 1, \quad r, \quad \varrho, \quad \varrho', \quad r\varrho, \quad r\varrho' \]
The cyclic permutation $(k_0, k_1, k_2, k_3)$ corresponds
to the permutation
\[
\begin{array}{lll}
r, & \varrho, & \varrho',\\
-r, & \varrho', & -\varrho,
\end{array}
\]
and if one repeats this permutation, one sees
that a cyclic function of the $k$ cannot change
if, in the linear representation by these six quantities,
the following permutations are made:
\[
\begin{array}{llllll}
1, & r, & \varrho, & \varrho', & r\varrho, & r\varrho'\\
1, & -r, & \varrho', & -\varrho, & -r\varrho', & r\varrho\\
1, & r, & -\varrho, & -\varrho', & -r\varrho, & -r\varrho'\\
1, & -r, & -\varrho', & \varrho, & r\varrho', & -r\varrho,
\end{array}
\]
and if one adds the four expressions thus obtained,
one obtains a rational expression for the cyclic function
of the $k$.

If one wishes to form the biquadratic equation whose roots
are the $k$, it is best to introduce $k - C = x$ as the unknown
%CONT>>
